\section{The valence formula}\label{sec:valence}

We now state the headline result of the modular side of the project:
the classical valence formula for weight-$k$ modular forms on
$\operatorname{SL}_2(\mathbb{Z})$. The argument follows the textbook
proof of Diamond--Shurman, Theorem 3.1.1: take the logarithmic
derivative $f'/f$ of $f$ around the fundamental-domain boundary at
height $H$, identify each residue with an order of vanishing, and
match the resulting weighted sum against $k/12$ via the
Hungerbühler--Wasem residue theorem and modular invariance.

\begin{theorem}[Orbit-sum identity from the PV chain]
  \label{thm:valence-orbit-sum-of-pvChain}
  \lean{valence_formula_orbit_sum_of_pvChain}
  \uses{def:generalized-winding-number, thm:gwn-at-i, thm:gwn-at-rho,
    thm:gwn-at-rho-plus-one}
  \leanok
  Let $k \in \mathbb{Z}$ and let $f$ be a weight-$k$ modular form on
  $\operatorname{SL}_2(\mathbb{Z})$ ($f \in M_k(\Gamma(1))$). Let
  $S \subseteq \mathcal{D}$ be a finite subset of the standard
  fundamental domain such that every order-non-zero point of $f$ in
  $\mathcal{D}$ lies in $S$. Assume there exists a height $H$ and a
  full winding-data bundle $D$ on $\partial \mathcal{D}_H$ (i.e.\ the
  interior winding number $-1$, the elliptic winding numbers
  $-\tfrac{1}{2}, -\tfrac{1}{6}, -\tfrac{1}{6}$ at $i, \rho, \rho+1$
  respectively, and the smooth-boundary winding $-\tfrac{1}{2}$ at
  every other on-curve point of $\partial \mathcal{D}_H$) for which
  every $s \in S$ has $\operatorname{Im}(s) < H$ and the
  PV chain identity
  \[
    \sum_{s \in S}
      \operatorname{gWN}\!\bigl(D.\mathrm{boundary},\, s\bigr) \cdot
      \operatorname{ord}_{s}(f)
    \;=\; -\bigl(k/12 - \operatorname{ord}_\infty(f)\bigr)
  \]
  holds. Then
  \[
    \operatorname{ord}_\infty(f) + \tfrac{1}{2}\,\operatorname{ord}_{i}(f)
      + \tfrac{1}{3}\,\operatorname{ord}_{\rho}(f)
      + \sum_{s \in \mathrm{INT}}\operatorname{ord}_{s}(f)
      + \sum_{s \in \mathrm{sLeftVert}} \operatorname{ord}_{s}(f)
      + \sum_{s \in \mathrm{LeftArc}} \operatorname{ord}_{s}(f)
    = \frac{k}{12},
  \]
  where
  \begin{itemize}
    \item $\mathrm{INT}$ is the subset of $S$ of strictly interior
      points (non-elliptic, $\|s\| > 1$, $|\operatorname{Re} s| < 1/2$);
    \item $\mathrm{sLeftVert}$ is the canonical-rep subset of $S$
      on the left vertical edge $\operatorname{Re} s = -1/2$;
    \item $\mathrm{LeftArc}$ is the subset of $S$ on the left half of
      the unit-circle arc $\|s\| = 1$, $\operatorname{Re} s < 0$ and
      $s \ne \rho$.
  \end{itemize}
\end{theorem}

\begin{proof}
  \uses{thm:gwn-at-i, thm:gwn-at-rho, thm:gwn-at-rho-plus-one,
    thm:hw-3-3-clean-full-mero}
  \leanok
  The PV chain identity gives a weighted sum of orders over $S$
  matching $-(k/12 - \operatorname{ord}_\infty(f))$. Substituting the
  three explicit elliptic winding numbers ($-\tfrac{1}{2}$ at $i$,
  $-\tfrac{1}{6}$ at $\rho$ and at $\rho+1$, see
  \cref{thm:gwn-at-i,thm:gwn-at-rho,thm:gwn-at-rho-plus-one}), and
  using the interior winding number $-1$ at strictly interior points
  and the smooth-boundary winding $-\tfrac{1}{2}$ at every other
  on-curve point, replaces $\operatorname{gWN}$ in each term by an
  explicit numerical weight. The right and left boundary subsets are
  paired via modular invariance (the $T$-translation identifies the
  two vertical edges, and the $S$-action identifies the two halves
  of the unit-circle arc), reducing each pair to a single canonical
  representative. The orders at $\rho$ and $\rho+1$ are equal by
  modular invariance; absorbing this identity gives the
  $\tfrac{1}{3}\operatorname{ord}_\rho(f)$ term. Rearranging yields
  the claimed equality.
\end{proof}

\begin{theorem}[Valence formula -- textbook unconditional form]
  \label{thm:valence-formula-textbook}
  \lean{valence_formula_textbook}
  \uses{thm:valence-orbit-sum-of-pvChain}
  \leanok
  Let $k \in \mathbb{Z}$ and let $f \in M_k(\Gamma(1))$ be a weight-$k$
  modular form on the full modular group $\operatorname{SL}_2(\mathbb{Z})$
  with $f \ne 0$. Then
  \[
    \operatorname{ord}_\infty(f) +
      \tfrac{1}{2}\,\operatorname{ord}_{i}(f) +
      \tfrac{1}{3}\,\operatorname{ord}_{\rho}(f) +
      \sum_{q \in \mathrm{NonEllOrbit}} \operatorname{ord}_{q}(f)
    \;=\; \frac{k}{12},
  \]
  where the last sum runs over the $\operatorname{SL}_2(\mathbb{Z})$-orbits
  in $\mathbb{H}$ that are neither the orbit of $i$ nor the orbit of
  $\rho$, and $\operatorname{ord}_q(f)$ is the order of vanishing of
  $f$ at any representative of $q$ (well-defined by modular
  invariance).
\end{theorem}

\begin{proof}
  \uses{thm:valence-orbit-sum-of-pvChain, thm:hw-3-3-clean-full-mero}
  \leanok
  The proof is a culmination of the entire Hungerbühler--Wasem
  apparatus and the orbit machinery. Apply the
  Hungerbühler--Wasem generalised residue theorem
  (\cref{thm:hw-3-3-clean-full-mero}) to the function $f'/f$ around
  the fundamental-domain boundary $\partial \mathcal{D}_H$ at a
  sufficiently large height $H$. The simple poles of $f'/f$ on
  $\mathcal{D}$ are precisely the zeros of $f$; the residue of $f'/f$
  at a zero of $f$ is the order of vanishing. The
  Hungerbühler--Wasem identity therefore takes the form
  \[
    \operatorname{PV}_S\!\!\int_{\partial \mathcal{D}_H}
      \frac{f'(z)}{f(z)}\,\mathrm{d}z
    \;=\; \sum_{s \in S} 2 \pi i \,
      \operatorname{gWN}(\partial \mathcal{D}_H, s) \,
        \operatorname{ord}_{s}(f).
  \]
  The left-hand side is in turn computed by the modular contour
  integral around the closed boundary $\partial \mathcal{D}_H$ to
  equal $2\pi i \, (\operatorname{ord}_\infty(f) - k/12)$ — the
  textbook contour-integral calculation, in which the top horizontal
  segment contributes $\operatorname{ord}_\infty(f)$ via the
  $q$-expansion, the two vertical edges cancel by $T$-invariance, and
  the two arcs are related by $S$-invariance, summing to $-k/12$.
  Equating the two sides and dividing by $2 \pi i$ gives the PV
  chain identity that hypothesises
  \cref{thm:valence-orbit-sum-of-pvChain}. The orbit-sum identity
  of that theorem rewrites the right-hand side as
  \[
    \operatorname{ord}_\infty(f) +
      \tfrac{1}{2}\operatorname{ord}_i(f) +
      \tfrac{1}{3}\operatorname{ord}_\rho(f) +
      \sum_{q \in \mathrm{NonEllOrbit}} \operatorname{ord}_{q}(f) = k/12,
  \]
  which is the textbook valence formula.
\end{proof}
